\subsection{Family Classification at a Common Alarm}
\label{sec:second-family-design}

We froze the history bank, the fixed-lag estimate
\eqref{eq:alarm-family-mass} and an adaptation of Li and Wang's CNN-F
before simulating 36 new two-event records on the second plant. Nine
records covered each ordered pair of load and line families. The first
intervention cleared 0.5--3\,s before the second began. Each record used
a different quiet startup from earlier runs on this plant, then fresh instrument
innovations and new draws of sites, switch times and load sizes. The
physical state evolved continuously through both interventions.

Li and Wang's classifier uses a 30-report window starting at a
trigger \cite{liwang2019successive}. All arms
received the history bank's causal alarms. At report $b+29$, we scored
the original family named at $b$, the fixed-lag family and CNN-F.
The common detector retained its calibration budget of one false alarm
per hour of quiet operation (one-sided 95\% upper bound).
We matched alarms by time alone and counted a missing, late or
unsupported diagnosis as an error. The deadline was the true clearance
of the second event. A second pair of arms used the same alarms,
publication times and fixed-lag rule. Both merged the recovery histories
at each reported restoration. One also zeroed the physical part of the
merged mean. This comparison varies the physical state supplied to the
next prediction while holding the classifier and the decision times fixed.

CNN-F received eight voltage magnitudes, eight frequencies and seven
active-power proxies constructed from the usable terminal phasors. It
extrapolated the previous response by dynamic mode decomposition, then
classified eigenvalue and singular-value features of the residual.
We trained the two convolutional paths on 360 independent, isolated
load and line interventions on the second plant and selected weights
on 108 further isolated interventions. We also ran the trained classifier
with subtraction disabled. Its two-class weights were fitted for these
PMUs. The published IEEE 68-bus weights have different class labels.
The spectrum here describes the evolution of the reported phasors.
The simulator exports fundamental-frequency phasors at 30 reports/s.
Different current waveforms can share those phasors and have different
harmonic distortion, so the comparison uses no current THD feature.
